
\documentclass{article}
\usepackage{colortbl}
\usepackage{makecell}
\usepackage{multirow}
\usepackage{supertabular}

\begin{document}

\newcounter{utterance}

\centering \large Interaction Transcript for game `codenames', experiment `association\_difficult', episode 2 with reflexion{-}llama{-}3.1{-}8b{-}instruct.
\vspace{24pt}

{ \footnotesize  \setcounter{utterance}{1}
\setlength{\tabcolsep}{0pt}
\begin{supertabular}{c@{$\;$}|p{.15\linewidth}@{}p{.15\linewidth}p{.15\linewidth}p{.15\linewidth}p{.15\linewidth}p{.15\linewidth}}
    \# & $\;$A & \multicolumn{4}{c}{Game Master} & $\;\:$B\\
    \hline

    \theutterance \stepcounter{utterance}  
    & & & \multicolumn{2}{p{0.3\linewidth}}{
        \cellcolor[rgb]{0.9,0.9,0.9}{
            \makecell[{{p{\linewidth}}}]{
                \texttt{\tiny{[GM$|$GM]}}
                {'hidden': {'team': ['jacks', 'opossum', 'window', 'meteorology', 'cello'], 'innocent': ['puzzle', 'top', 'actor', 'cougar'], 'opponent': [], 'assassin': ['hat']}, 'revealed': {'team': {'team': [], 'innocent': [], 'opponent': [], 'assassin': []}, 'opponent': {'team': [], 'innocent': [], 'opponent': [], 'assassin': []}}}
            }
        }
    }
    & & \\ \\

    \theutterance \stepcounter{utterance}  
    & & \multicolumn{4}{p{0.6\linewidth}}{
        \cellcolor[rgb]{0.9,0.9,0.9}{
            \makecell[{{p{\linewidth}}}]{
                \texttt{\tiny{[P1$\langle$GM]}}
                \texttt{You have attempted this same task before and did not succeed. These are the plans you wrote after those attempts.} \\
\texttt{Trial \#1: New plan: } \\
\\ 
\texttt{In my previous attempts, I focused on providing clues that were directly related to the words on the list, but I didn't consider the potential consequences of giving away too much information. For example, when I gave the clue "instrument" and related it to multiple words, my teammate was able to guess the word "violin" which is an innocent word, but it ended my teammate's turn. Similarly, when I gave the clue "academic" and related it to the words "professor" and "meteorology", my teammate was able to guess the word "professor" correctly, but it also ended my teammate's turn.} \\
\\ 
\texttt{Next time, I will focus on providing clues that are more general and less likely to give away too much information. I will also try to provide clues that are more likely to lead to the opponent's words, rather than the innocent words. For example, I might give a clue that is related to a category or a theme that is more likely to include the opponent's words. I will also try to pay closer attention to the words that have been removed from the lists and avoid giving clues that are too similar to the words that have already been guessed.} \\
\\ 
\\ 
\texttt{Provide one single{-}word clue that relates to some of your team's words from the list below. You can choose to how many and to which words your clue relates to.} \\
\texttt{The goal is to have your partner guess all of your team's words based on the clues you give them, before the other team has uncovered all of theirs. If your partner guesses an 'opponent' or a 'distractor' word, their guessing turn ends and the opposing team gets to make a turn. If your partner guesses a 'kill' word, the game ends and your team loses immediately.} \\
\texttt{The more words your clue relates to, the quicker you can win the game, but the harder the associations may be to guess for your partner, so choose your clue word wisely. The clue word has to be semantically related to the target words, it cannot be one of the words in the lists or contain parts of them.} \\
\\ 
\texttt{Always give your single{-}word clue and your comma{-}separated list of related target words in the following format and make your answers as short as possible, never include any other text than is required in this form:} \\
\\ 
\texttt{CLUE: <WORD>} \\
\texttt{TARGETS: <TARGETS>} \\
\\ 
\texttt{Your team words are: jacks, opossum, window, meteorology, cello.} \\
\texttt{Your opponent's team words are: .} \\
\texttt{Distractor words are: puzzle, top, actor, cougar.} \\
\texttt{The kill words are: hat.} \\
            }
        }
    }
    & \\ \\

    \theutterance \stepcounter{utterance}  
    & \multicolumn{4}{p{0.6\linewidth}}{
        \cellcolor[rgb]{0.9,0.9,0.9}{
            \makecell[{{p{\linewidth}}}]{
                \texttt{\tiny{[P1$\rangle$GM]}}
                \texttt{CLUE: animal} \\
\texttt{TARGETS: opossum, jacks} \\
            }
        }
    }
    & & \\ \\

    \theutterance \stepcounter{utterance}  
    & & & \multicolumn{2}{p{0.3\linewidth}}{
        \cellcolor[rgb]{0.9,0.9,0.9}{
            \makecell[{{p{\linewidth}}}]{
                \texttt{\tiny{[GM$|$GM]}}
                \texttt{animal} \\
            }
        }
    }
    & & \\ \\

    \theutterance \stepcounter{utterance}  
    & & & \multicolumn{2}{p{0.3\linewidth}}{
        \cellcolor[rgb]{0.9,0.9,0.9}{
            \makecell[{{p{\linewidth}}}]{
                \texttt{\tiny{[GM$|$GM]}}
                ['opossum', 'jacks']
            }
        }
    }
    & & \\ \\

    \theutterance \stepcounter{utterance}  
    & & & \multicolumn{2}{p{0.3\linewidth}}{
        \cellcolor[rgb]{0.9,0.9,0.9}{
            \makecell[{{p{\linewidth}}}]{
                \texttt{\tiny{[GM$|$GM]}}
                {'word': 'opossum', 'assignment': 'team'}
            }
        }
    }
    & & \\ \\

    \theutterance \stepcounter{utterance}  
    & & & \multicolumn{2}{p{0.3\linewidth}}{
        \cellcolor[rgb]{0.9,0.9,0.9}{
            \makecell[{{p{\linewidth}}}]{
                \texttt{\tiny{[GM$|$GM]}}
                {'word': 'jacks', 'assignment': 'team'}
            }
        }
    }
    & & \\ \\

    \theutterance \stepcounter{utterance}  
    & & & \multicolumn{2}{p{0.3\linewidth}}{
        \cellcolor[rgb]{0.9,0.9,0.9}{
            \makecell[{{p{\linewidth}}}]{
                \texttt{\tiny{[GM$|$GM]}}
                \texttt{opponent has won} \\
            }
        }
    }
    & & \\ \\

    \theutterance \stepcounter{utterance}  
    & \multicolumn{4}{p{0.6\linewidth}}{
        \cellcolor[rgb]{0.95,0.95,0.95}{
            \makecell[{{p{\linewidth}}}]{
                \texttt{\tiny{[P2$\rightarrow$P1]}}
                \texttt{Post{-}episode Reflexion (not sent to GM):} \\
\texttt{New plan: } \\
\\ 
\texttt{In my previous attempts, I focused on providing clues that were directly related to the words on the list, but I didn't consider the potential consequences of giving away too much information. For example, when I gave the clue "animal" and related it to the words "opossum" and "jacks", I didn't think about the fact that "animal" is a very general term that could relate to many words on the list, including the opponent's words.} \\
\\ 
\texttt{Next time, I will focus on providing clues that are more specific and less likely to give away too much information. I will also try to provide clues that are more likely to lead to the opponent's words, rather than the innocent words. For example, I might give a clue that is related to a category or a theme that is more likely to include the opponent's words. I will also try to pay closer attention to the words that have been removed from the lists and avoid giving clues that are too similar to the words that have already been guessed.} \\
\\ 
\texttt{To achieve this, I will start by looking for common themes or categories that relate to multiple words on the list, including the opponent's words. I will then try to come up with a clue that is specific to that category or theme, but not too specific that it gives away too much information. I will also try to avoid giving clues that are too similar to the words that have already been guessed, and pay closer attention to the words that have been removed from the lists.} \\
\\ 
\texttt{For example, if I notice that the opponent's words are related to a specific category, such as "animals that live in trees", I might give a clue that is related to that category, such as "arboreal". This way, I can provide a clue that is specific to the category, but not too specific that it gives away too much information.} \\
            }
        }
    }
    & & \\ \\

\end{supertabular}
}

\end{document}
